% ═══════════════════════════════════════════════════════════════════════════
% AB: Repaso Unidad 1 + 2 (Backup für LP)
% ═══════════════════════════════════════════════════════════════════════════

\documentclass[a4paper,11pt]{article}

\usepackage[utf8]{inputenc}
\usepackage[T1]{fontenc}
\usepackage[ngerman]{babel}
\usepackage{helvet}
\renewcommand{\familydefault}{\sfdefault}

\usepackage[margin=18mm]{geometry}
\usepackage{xcolor}
\usepackage{tikz}
\usepackage{enumitem}
\usepackage{array}
\usepackage{booktabs}
\usepackage{multicol}
\usepackage{tabularx}

% Fachfarbe Spanisch
\definecolor{spanisch}{HTML}{c41e3a}
\definecolor{hellgrau}{HTML}{f5f5f5}
\definecolor{mittelgrau}{HTML}{888888}

% Lücke
\newcommand{\luecke}[1]{\underline{\hspace{#1}}}

% Punktebox
\newcommand{\punktebox}[1]{\hfill\fbox{\small\textcolor{mittelgrau}{\_\_\_\,/\,#1}}}

% Aufgaben-Umgebung
\newenvironment{aufgabe}[2]{%
    \vspace{3mm}
    \noindent\textbf{\color{spanisch}Aufgabe #1: #2}
    \vspace{1mm}
    \par
}{\vspace{2mm}}

\setlength{\parindent}{0pt}
\setlength{\parskip}{0pt}

\pagestyle{empty}

\begin{document}

% ═══════════════════════════════════════════════════════════════════════════
% KOPF
% ═══════════════════════════════════════════════════════════════════════════

\noindent
\begin{tabularx}{\textwidth}{Xr}
    {\Large\bfseries\color{spanisch}Repaso: Unidad 1 + 2} & {\small Klasse 9 -- Spanisch} \\
\end{tabularx}

\vspace{1mm}
\noindent\rule{\textwidth}{0.4pt}
\vspace{1mm}

\noindent
{\small Nombre: \luecke{5cm} \hfill Fecha: \luecke{3cm} \hfill Punkte: \luecke{1cm} / 50}

\vspace{4mm}

% ═══════════════════════════════════════════════════════════════════════════
% AUFGABE 1: Vokabeln Restaurante
% ═══════════════════════════════════════════════════════════════════════════

\begin{aufgabe}{1}{Vocabulario -- El restaurante \punktebox{5}}

Verbinde die spanischen Wörter mit der deutschen Bedeutung.

\vspace{2mm}
\begin{multicols}{2}
\begin{tabular}{rl@{\hspace{8mm}}rl}
1. & el camarero & a) & die Gabel \\[1mm]
2. & la cuenta & b) & der Kellner \\[1mm]
3. & el postre & c) & das Trinkgeld \\[1mm]
4. & el tenedor & d) & die Rechnung \\[1mm]
5. & la propina & e) & der Nachtisch \\
\end{tabular}

\vspace{3mm}
\textbf{Lösung:} 1\luecke{0.5cm} 2\luecke{0.5cm} 3\luecke{0.5cm} 4\luecke{0.5cm} 5\luecke{0.5cm}
\end{multicols}
\end{aufgabe}

% ═══════════════════════════════════════════════════════════════════════════
% AUFGABE 2: Vokabeln Casa
% ═══════════════════════════════════════════════════════════════════════════

\begin{aufgabe}{2}{Vocabulario -- La casa y los muebles \punktebox{6}}

Übersetze ins Deutsche.

\vspace{2mm}
\begin{multicols}{2}
\begin{enumerate}[itemsep=1mm]
    \item el salón = \luecke{3cm}
    \item la cocina = \luecke{3cm}
    \item el armario = \luecke{3cm}
    \item la nevera = \luecke{3cm}
    \item el espejo = \luecke{3cm}
    \item la ducha = \luecke{3cm}
\end{enumerate}
\end{multicols}
\end{aufgabe}

% ═══════════════════════════════════════════════════════════════════════════
% AUFGABE 3: pedir / servir
% ═══════════════════════════════════════════════════════════════════════════

\begin{aufgabe}{3}{Verbos: pedir y servir (e $\rightarrow$ i) \punktebox{6}}

Ergänze die richtige Verbform.

\vspace{2mm}
\begin{enumerate}[itemsep=2mm]
    \item Yo \luecke{2.5cm} una paella. (pedir)
    \item El camarero \luecke{2.5cm} la comida. (servir)
    \item Nosotros \luecke{2.5cm} la cuenta. (pedir)
    \item Tú \luecke{2.5cm} un café. (pedir)
    \item Ellos \luecke{2.5cm} en el restaurante. (servir)
    \item Vosotros \luecke{2.5cm} muy bien. (servir)
\end{enumerate}
\end{aufgabe}

% ═══════════════════════════════════════════════════════════════════════════
% AUFGABE 4: conocer / parecer
% ═══════════════════════════════════════════════════════════════════════════

\begin{aufgabe}{4}{Verbos: conocer y parecer (-zco) \punktebox{6}}

Ergänze die richtige Verbform.

\vspace{2mm}
\begin{enumerate}[itemsep=2mm]
    \item Yo \luecke{2.5cm} a María. (conocer)
    \item ¿Qué te \luecke{2.5cm} la película? (parecer)
    \item Nosotros \luecke{2.5cm} bien Madrid. (conocer)
    \item Me \luecke{2.5cm} interesantes los libros. (parecer)
    \item Tú no \luecke{2.5cm} a mi hermano. (conocer)
    \item Yo \luecke{2.5cm} simpático. (parecer)
\end{enumerate}
\end{aufgabe}

% ═══════════════════════════════════════════════════════════════════════════
% AUFGABE 5: Pretérito perfecto regulär
% ═══════════════════════════════════════════════════════════════════════════

\begin{aufgabe}{5}{El pretérito perfecto -- regulär \punktebox{5}}

Bilde das pretérito perfecto (haber + Partizip).

\vspace{2mm}
\begin{enumerate}[itemsep=2mm]
    \item Hoy yo \luecke{3.5cm} mucho. (trabajar)
    \item Esta semana nosotros \luecke{3.5cm} en el restaurante. (comer)
    \item ¿Tú ya \luecke{3.5cm} la tarea? (terminar)
    \item Ellos nunca \luecke{3.5cm} a España. (viajar)
    \item María \luecke{3.5cm} un libro interesante. (leer)
\end{enumerate}
\end{aufgabe}

% ═══════════════════════════════════════════════════════════════════════════
% AUFGABE 6: Pretérito perfecto irregulär
% ═══════════════════════════════════════════════════════════════════════════

\begin{aufgabe}{6}{El pretérito perfecto -- irregulär \punktebox{5}}

Bilde das pretérito perfecto mit unregelmäßigen Partizipien.

\vspace{2mm}
\begin{enumerate}[itemsep=2mm]
    \item ¿Qué \luecke{3.5cm} hoy? (hacer, tú)
    \item Nosotros \luecke{3.5cm} la verdad. (decir)
    \item Yo no \luecke{3.5cm} la película. (ver)
    \item Ella \luecke{3.5cm} una carta. (escribir)
    \item ¿Dónde \luecke{3.5cm} las llaves? (poner, tú)
\end{enumerate}
\end{aufgabe}

\vspace{2mm}
\noindent\fcolorbox{mittelgrau}{hellgrau}{\parbox{0.97\textwidth}{\small\textbf{Hilfe:} hacer $\rightarrow$ hecho, decir $\rightarrow$ dicho, ver $\rightarrow$ visto, escribir $\rightarrow$ escrito, poner $\rightarrow$ puesto}}

% ═══════════════════════════════════════════════════════════════════════════
% AUFGABE 7: Direktobjektpronomen
% ═══════════════════════════════════════════════════════════════════════════

\begin{aufgabe}{7}{Los pronombres de objeto directo (lo/la/los/las) \punktebox{8}}

Ersetze das Objekt durch das passende Pronomen.

\vspace{2mm}
\begin{enumerate}[itemsep=2mm]
    \item ¿El libro? \luecke{1.5cm} leo todos los días.
    \item ¿La mesa? \luecke{1.5cm} ponemos en el salón.
    \item ¿El sofá? \luecke{1.5cm} compramos ayer.
    \item ¿La lámpara? \luecke{1.5cm} pongo en el dormitorio.
    \item ¿Los muebles? \luecke{1.5cm} subimos al piso.
    \item ¿Las sillas? \luecke{1.5cm} ponemos en la cocina.
    \item ¿Los cuadros? \luecke{1.5cm} colgamos en el salón.
    \item ¿Las plantas? \luecke{1.5cm} riego todos los días.
\end{enumerate}
\end{aufgabe}

% ═══════════════════════════════════════════════════════════════════════════
% AUFGABE 8: Bewegungsverben
% ═══════════════════════════════════════════════════════════════════════════

\begin{aufgabe}{8}{Los verbos de movimiento + preposiciones \punktebox{9}}

Ergänze \textbf{de}, \textbf{del}, \textbf{a}, \textbf{al} oder \textbf{en}.

\vspace{2mm}
\begin{enumerate}[itemsep=2mm]
    \item Bajo \luecke{1.5cm} coche. (el coche)
    \item Subo \luecke{1.5cm} piso. (el piso)
    \item Saco los libros \luecke{1.5cm} armario. (el armario)
    \item Llevo la mesa \luecke{1.5cm} salón. (el salón)
    \item Pongo el libro \luecke{1.5cm} la estantería.
    \item La lámpara está \luecke{1.5cm} el dormitorio.
    \item Traigo la silla \luecke{1.5cm} la cocina. (von der Küche)
    \item Ponemos el sofá \luecke{1.5cm} el salón.
    \item Bajamos las cajas \luecke{1.5cm} sótano. (el sótano = Keller)
\end{enumerate}
\end{aufgabe}

\vspace{3mm}
\noindent\fcolorbox{mittelgrau}{hellgrau}{\parbox{0.97\textwidth}{\small\textbf{Hilfe:} \textbf{de/del} = weg von | \textbf{a/al} = hin zu | \textbf{en} = in (Position)}}

\vfill

\noindent\rule{\textwidth}{0.4pt}
{\small\color{mittelgrau} Backup-AB für LP-repaso-U1-U2 | Klasse 9 | Spanisch}

\end{document}
